\Needspace{18\baselineskip}
\section{Concrete graph examples and regression proofs}
\Needspace{24\baselineskip}
\subsection{Self-loop gluing and the necessary reflection restriction}
Use a singleton wire-label set and a singleton operation alphabet, whose only
operation has one input and one output of that label. The graph named
$\mathsf{loop}$ has one internal vertex, one edge, no outer boundary, and that
vertex as both the edge's source and target.
\lean{Graph/Examples}{7}{9}{Examples.unary}
\lean{Graph/Examples}{12}{44}{Examples.loop}
Its finite carriers are both $\operatorname{Fin}(1)$. Singleton incidence lists
are distinct; all edges are equal since their indices are bounded by one. The
only internal vertex is incident on both sides, providing coverage, while
boundary cases have empty finite-index types. Label typing is reflexive. Thus
this is a well-formed net despite its cycle. Let $\mathsf{wireEdge}$ be the
corolla of the same unary operation, with its two distinct boundary vertices.
\lean{Graph/Examples}{46}{46}{Examples.wireEdge}

The corolla is not adjacent to itself: a shared target/source witness would have
to be both its output constructor and its input constructor, an impossible
constructor equality.
\lean{Graph/Examples}{49}{54}{Examples.corolla_not_self_adjacent}
When the corolla replaces the loop's sole edge, its input and output images are
equal. Both are, by definition, the same old internal vertex, because the host
hole's source and target port coincide. Hence the inserted edge is adjacent to
itself: use that old vertex as the witness and compute the mapped singleton
source and target lists.
\lean{Graph/Examples}{58}{60}{Examples.inserted_boundary_identified}
\lean{Graph/Examples}{62}{67}{Examples.inserted_self_adjacent}
Consequently the inserted two-entry list $[e,e]$ is a path, while $[e,e]$ is not
a path in the original corolla. The positive part uses its new self-adjacency and
a singleton tail path. The negative part extracts its first adjacency and
contradicts the corolla result. This proves that unconditional reflection of
inserted paths would be false; the non-self-loop hypothesis cannot simply be
omitted from that theorem.
\lean{Graph/Examples}{71}{76}{Examples.unconditional_path_reflection_fails}
The original loop itself has the two-entry path $[e,e]$, using its sole internal
vertex as the shared witness. This verifies that the graph representation has
not silently imposed global acyclicity.
\lean{Graph/Examples}{80}{81}{Examples.loop_is_cyclic}

The finite ordinal with one element is bijective to the unit type: send its
index to the sole unit element and send unit back to zero. The inverse laws
follow respectively from the ordinal bound and unit case analysis. An anonymous
example reindexes the corolla's edge carrier along this bijection and applies
shallow plugging congruence. It obtains isomorphic filled loops despite the
different replacement-edge carrier types, testing actual renaming rather than
reflexivity alone.
\lean{Graph/Examples}{83}{87}{Examples.oneToUnit}
\lean{Graph/Examples}{89}{91}{Examples.example (line 89)}

\Needspace{24\baselineskip}
\subsection{Values, boxes, and distinguished occurrence addresses}
For the linear examples, take two operation labels, value and other-value. Give
both zero term arguments and no thunk arguments; each therefore has signature
$*\to []$. The graph $\mathsf{value}$ is the corolla of the first operation at
depth zero.
\lean{Graph/Examples}{99}{100}{Linear.Examples.Op}
\lean{Graph/Examples}{102}{104}{Linear.Examples.operations}
\lean{Graph/Examples}{106}{107}{Linear.Examples.value}
Boxing it with zero bound and free variables yields a closed box of type
$T^0(*)\to []$ at depth one. A further example constructs the analogous graph
at depth bound two by putting a primitive value at depth one inside a box.
This is a larger \emph{available bound}, not a claim that the example has two
actually nested boxes.
\lean{Graph/Examples}{110}{110}{Linear.Examples.boxedValue}
\lean{Graph/Examples}{112}{115}{Linear.Examples.higherBoxedValue}
An anonymous proof verifies that any occurrence in the depth-one boxed value
belongs to the recursive occurrence enumeration, by the general completeness
theorem. The outer box edge is represented by a shallow left-tagged occurrence;
the value edge in its body is a right-tagged pair of parent edge and child edge.
\lean{Graph/Examples}{117}{118}{Linear.Examples.example (line 117)}
\lean{Graph/Examples}{122}{122}{Linear.Examples.outerOccurrence}
\lean{Graph/Examples}{124}{124}{Linear.Examples.innerOccurrence}
These two addresses are unequal. If equal, applying the sum-tag test would give
$\mathsf{true}=\mathsf{false}$, which is impossible. This checks separation of
outer and inner occurrences even though each individual edge carrier is singleton.
\lean{Graph/Examples}{126}{129}{Linear.Examples.outer_ne_inner}

The other-value corolla has exactly the same interface as the value corolla.
A site inside the boxed value descends through its only box edge and selects the
value edge in its body. Replacing at this site with other-value defines
$\mathsf{replacedInside}$.
\lean{Graph/Examples}{131}{132}{Linear.Examples.otherValue}
\lean{Graph/Examples}{135}{136}{Linear.Examples.insideValue}
\lean{Graph/Examples}{138}{139}{Linear.Examples.replacedInside}
Its outer source and target functions equal those of the original boxed value;
instantiate the two deep-locality theorems at the selected site. Its sole outer
edge label is the box with body $\mathsf{value}[e:=\mathsf{otherValue}]$.
For this label equation, unfold the site and replacement, simplify the equality
test selecting the sole outer edge, and use reflexivity for the resulting body.
The outer box is retained rather than replaced by the value operation itself.
\lean{Graph/Examples}{141}{142}{Linear.Examples.deep_replacement_source_unchanged}
\lean{Graph/Examples}{144}{145}{Linear.Examples.deep_replacement_target_unchanged}
\lean{Graph/Examples}{147}{151}{Linear.Examples.deep_replacement_body}

\Needspace{24\baselineskip}
\subsection{An input-rooted operation path}
Give the single unit operation one term argument and no binders. Its corolla
$\mathsf{argument}$ has type $*\to *$. The input, its only edge, and the output
form a forward operation path: read this as the one-step reverse trace, check
membership in the two singleton incidence lists, and end with an empty trace.
\lean{Graph/Examples}{153}{155}{Linear.Examples.argumentOperations}
\lean{Graph/Examples}{157}{158}{Linear.Examples.argument}
\lean{Graph/Examples}{160}{166}{Linear.Examples.argument_path}
Any reverse operation trace from its output back to that output has empty edge
list. Apply the no-cycle theorem to the constructed input-rooted path, deriving
the required unary condition from the typed signature. This exercises the
connection between concrete graph typing and generic machine terminal uniqueness.
\lean{Graph/Examples}{168}{172}{Linear.Examples.no_cycle_after_argument}

\Needspace{24\baselineskip}
\subsection{Filling a hole inside a certified context}
Choose the typed hole label with name zero, one term input, and no outputs.
Its corolla $\mathsf{singleHole}$ is a depth-zero hypernet. It has unique holes
because any two edge occurrences are equal: both are elements of
$\operatorname{Fin}(1)$, so their underlying indices are zero. Pairing this graph
with the proof constructs a context.
\lean{Graph/Examples}{174}{174}{Linear.Examples.holeLabel}
\lean{Graph/Examples}{176}{177}{Linear.Examples.singleHole}
\lean{Graph/Examples}{179}{184}{Linear.Examples.singleHole_unique}
\lean{Graph/Examples}{186}{186}{Linear.Examples.singleContext}
Box the graph with zero bound and free variables to obtain $\mathsf{boxedHole}$,
and box the certified context to obtain $\mathsf{boxedContext}$. The unique-hole
certificate in the latter is supplied by the general boxing theorem. Looking
up the deep address consisting of the only outer edge and the only body edge
returns the chosen hole label by definition.
\lean{Graph/Examples}{188}{189}{Linear.Examples.boxedHole}
\lean{Graph/Examples}{191}{192}{Linear.Examples.boxedContext}
\lean{Graph/Examples}{194}{195}{Linear.Examples.hole_is_deep}
The corresponding site descends into that box and chooses its shallow hole edge.
Filling this site with the value corolla produces $\mathsf{filledBox}$. Its outer
label is the same zero-bound, zero-free box, now containing
$\mathsf{singleHole}[e:=\mathsf{value}]$; unfold deep replacement and simplify the
selected-edge equality test to prove this equation.
\lean{Graph/Examples}{197}{198}{Linear.Examples.boxedHoleSite}
\lean{Graph/Examples}{200}{200}{Linear.Examples.filledBox}
\lean{Graph/Examples}{202}{206}{Linear.Examples.filled_box_body}
Finally every edge of this filled body has no primitive hole label. Split on
its edge tag. An old surviving edge would be an element of the singleton old
carrier different from its only element, contradicting the surviving-edge
certificate. An inserted value edge is an operation, whose primitive hole lookup
is definitionally empty. Thus this example proves removal of the old hole, not
merely the well-typedness of the filled graph.
\lean{Graph/Examples}{210}{216}{Linear.Examples.filled_body_no_holes}
